\section{Case Studies}\label{sec:usecases}

In this section we conclude our case studies. See \cref{sec:eval} for the setup
and methodology behind each study.


% ===========================================================
\subsection{``Battery Saver'' Mode}\label{subsec:usecase-battery}

% About: background
In this first case study, we revisit the motivational Forward switching example
(\cref{sec:motivation}). Our goal is to complete a workload within a strict
energy budget. Here, our netbook's storage will use {\em Freestyle Balanced
(FB)} when operating normally and {\em ChaCha8 (C8)} when battery saver mode is
activated.

\def \hmina {\hspace{-0.1in}}
\def \hminb {\hspace{-0.2in}}

\def \fgw {2in}
\def \fgh {1in}

% \begin{floatingfigure}[r]{2in}  (and \end{..})

\begin{figure}[t]
    \centerline{
        % \hmina
        % \includegraphics[width=\fgw]{data/usecase-battery}
        % \hminb
        \includegraphics[height=\fgh]{empty.eps}
    }

    \mycaption{fig:usecase-battery}{Battery saver case}{Energy-security
    tradeoff given strict energy budget as discussed in
    \cref{subsec:usecase-battery}.}
\end{figure}



% About: setup
To simulate this, we run the following twice: Write 10 40MB files using FB,
after 5s signal low battery. The workload is completed using homogeneous FDE on
the first run. On the second, the system enters battery saving mode and \sys
switches to C8 using the Forward switching model. \Cref{fig:usecase-battery}
shows the total energy used after completing both runs. {\em Homogeneous} FDE
(left) favors security even when backups are paused, while \sys (right) matches
the security and energy tradeoff to the situation. \uline{\sys results in a 3.1x
total energy reduction compared to prior work (homogeneous FDE), allowing us to
remain safely within our energy budget.}


% ===============================
\subsection{No-Downtime Encryption Upgrade}\label{subsec:usecase-upgrade}

% About: background
We revisit the Mirrored switching example (\cref{subsec:des-switch}). Our goal
is to upgrade live storage from one cipher to another without degraded service.
Our operator wants to upgrade from {\em ChaCha20 (C20)} to {\em Freestyle Fast
(FF)}.

% About: setup
To simulate pre- and post-migration state, we execute 10 40MB write-read
operations to two \sys instances using C20 and FF respectively. To simulate the
activity during migration, we repeat the operations on a third instance using
Mirrored switching.

\def \hmina {\hspace{-0.1in}}
\def \hminb {\hspace{-0.2in}}

\def \fgw {2in}
\def \fgh {1in}

% \begin{floatingfigure}[r]{2in}  (and \end{..})

\begin{figure}[t]
    \centerline{
        % \hmina
        % \includegraphics[width=\fgw]{usecase-upgrade}
        % \hminb
        \includegraphics[height=\fgh]{empty.eps}
    }

    \mycaption{fig:usecase-upgrade}{Cipher upgrade without downtime}{Switching
    storage system encryption with zero downtime as discussed in
    \cref{subsec:usecase-upgrade}.}
\end{figure}


% About: outcome
\Cref{fig:usecase-upgrade} shows I/O performance at the three stages of
migration. On the left we see similar read latency between the first two stages
since this model does not switch crypts for read I/Os until the migration
completes, after which performance matches the slower FF. On the right, we see
write latency increase to 8.1s during the migration as writes are duplicated;
post-migration, performance recovers to 6.1s. Thus, unlike inflexible prior
work, \sys \uline{successfully allows one to transition storage to another
cipher without interruption or egregious performance loss.}


% ================================================================
\subsection{Scalable Encryption}\label{subsec:usecase-scalable}

% About: background
We return to the Selective switching example
(\cref{subsec:des-switch}). Our goal is to achieve a performance win by ensuring
that the only data encrypted using the strongest, highest-overhead crypt is the
small percentage of data that requires such high security. We note the Selective
switching model requires users to annotate certain files with a special tag via
file system calls and stored in the inode. Because the \sys layer is
file-oblivious, every block I/O through the \sys layer will be labeled with the
corresponding crypt.

% About: setup
We begin with 10 5MB and 10 4KB write-read operations to \sys using {\em
Freestyle Strong (FS)}, a high security crypt. We repeat these operations on a
second instance using Selective switching where 30\% of the data is considered
highly sensitive and so uses FS while the rest uses {\em ChaCha8 (C8)}, a high
performance crypt.

\def \hmina {\hspace{-0.1in}}
\def \hminb {\hspace{-0.2in}}

\def \fgw {2in}
\def \fgh {1in}

% \begin{floatingfigure}[r]{2in}  (and \end{..})

\begin{figure}[t]
    \centerline{
        % \hmina
        {\begin{tikzpicture}[baseline]

    \pgfmathsetmacro{\ymax}{20} % set the maximum y value
    \pgfmathsetmacro{\ymaxbreak}{20.1} % set the y value at which overflow is drawn

    \begin{axis}[
        %axis x line*=bottom,
        height=1in,
        width=2.5in,
        tick align=outside,
        tick pos=bottom, % make sure ticks only appear at the bottom and left axes
        tick style={ black },
        y tick label style={ /pgf/number format/fixed, /pgf/number format/precision=0 },
        grid style={ dotted, gray },
        every node near coord/.append style={font=\tiny},
        %
        %
        % magic to make the numbers appear above the overly long bars:
        visualization depends on={rawy \as \rawy}, % save original y values
        % restrict y to domain*={ % now clip/restrict any y value to ymax
        %     \pgfkeysvalueof{/pgfplots/ymin}:\ymaxbreak
        % },
        % after end axis/.code={ % draw squiggly line indicating break
        %     \draw [semithick, white, decoration={snake,amplitude=0.1mm,segment length=0.75mm,post length=0.375mm}, decorate] (rel axis cs:0,1.01) -- (rel axis cs:1,1.01);
        % },
        % nodes near coords={\color{.!75!black}\pgfmathprintnumber\rawy}, % print the original y values (darkened in case they are too light)...
        % nodes near coords greater equal only=\ymax, % ... but ONLY if they are >= ymax
        % clip=false, % allow clip to protrude beyond ymax
        %
        every node near coord/.append style={font=\tiny,yshift=-0.5mm},
        nodes near coords={\color{black}\pgfmathprintnumber\rawy},
        nodes near coords align={vertical},
        %
        % Custom stuff to edit per template
        %
        %xlabel={\scriptsize Cipher Configuration},
        xlabel near ticks,
        xmin=FS, xmax=C8+FS,
        xtick=data,
        symbolic x coords={FS,C8+FS},
        xticklabels={Homogeneous,\sys},
        enlarge x limits=0.5, % add some breathing room along the x axis's sides
        %bar width=5mm,
        %major x tick style=transparent,
        %
        ylabel={\scriptsize Latency (s)},
        ylabel near ticks,
        %ylabel shift={-1mm},
        ylabel style={at={(ticklabel cs:0.4)}},
        ymajorgrids=true,
        ymin=0, ymax=30,
        ybar, % value will shift bars
        ytick={ 5, 15, 25 },
        %yticklabels={ 5, 15, 25},
        % extra y ticks={1},
        % extra y tick style={grid=major, grid style={dashed, black}},
        % extra y tick label={\empty},
        %bar width=4.5pt, % change size of bars
        %
        legend cell align=center,
        legend style={ column sep=1ex },
        legend entries={
            {\scriptsize 4K/reads},
            {\scriptsize 4K/writes},
            {\scriptsize 5M/reads},
            {\scriptsize 5M/writes},
        },
        legend style={
            draw=none,
            legend columns=4,
            at={(0.5,1.02)},
            anchor=south,
        },
    ]
        \addplot[fill=pinkDark, every node near coord/.append style={color=pinkDark}]
        table[x=conf, y=latr-4k, col sep=space] {data/usecase-scalable.dat};
        \addplot[fill=pinkDark, postaction={pattern=north east lines}, every node near coord/.append style={color=purpleDark}]
        table[x=conf, y=latw-4k, col sep=space] {data/usecase-scalable.dat};
        \addplot[fill=purpleDark, every node near coord/.append style={color=pinkDark}]
        table[x=conf, y=latr-5m, col sep=space] {data/usecase-scalable.dat};
        \addplot[fill=purpleDark, postaction={pattern=north east lines}, every node near coord/.append style={color=purpleDark}]
        table[x=conf, y=latw-5m, col sep=space] {data/usecase-scalable.dat};
    \end{axis}%
\end{tikzpicture}%
}
        % \hminb
        %\includegraphics[height=\fgh]{empty.eps}
    }

    \mycaption{fig:usecase-scalable}{Scalable encryption}{Filesystem-agnostic
    file-level encryption as discussed in \cref{subsec:usecase-scalable}.}
\end{figure}


% About: outcome
\Cref{fig:usecase-scalable} shows {\em Homogeneous} FDE (left) exhibits 5s to
20s latency depending on the workload while \sys (right) \uline{outperforms with
a reduction of 3.1x to 4.8x for read latency and 1.6x to 2.8x for write latency,
all without compromising on security}. Thus, \sys keeps our sensitive data at
the mandated security level while keeping the performance benefits of fast
crypto.
